\begin{solapartat}
És important que les línies de camp indiquin el sentit \punts{0,2} i que les superfícies equipotencials indiquin els valors dels seus potencials. \punts{0,2} No és necessari que el 0 correspongui a l'elèctrode negatiu.
% A l'original aquesta segona puntuació apareix entre claudàtors, «[0.2]», i no en requadre.

\figura[0.55]{sol_fig1.png}
\punts{0,2}

El valor del camp serà:
\[ E = \frac{\Delta V}{x} = \frac{300V}{0,2m} = 1,50\times 10^{3}\mathrm{V/m}\ \text{ó}\ 1,50\times 10^{3}\mathrm{N/C} \ \text{\punts{0,2}} \]
Les partícules negatives dipositades es mouran cap al pol positiu i les positives cap al pol negatiu. \punts{0,2}
\end{solapartat}

\begin{solapartat}
La força elèctrica ha de ser: $\overrightarrow{F} = q\overrightarrow{E} = -1,6\times 10^{-19}C\ 1500N/C\,\overrightarrow{i} = -2,40\times 10^{-16}N\,\overrightarrow{i}$ o bé $2,40\times 10^{-16}N\,\overrightarrow{i}$ si el signe de la càrrega és positiu. \punts{0,5}

Com que es mou amb un moviment rectilini i uniforme $\Rightarrow \Sigma\overrightarrow{F} = 0$, per tant la força de fricció ha de ser igual i de sentit contrari a la força elèctrica, o sigui el seu modul val: $2,40\times 10^{-16}N$ \punts{0,5}.
\end{solapartat}
